\documentclass[12pt,a4paper]{article}

\usepackage[utf8]{inputenc}
\usepackage[T1]{fontenc}
\usepackage[brazil]{babel}
\usepackage{geometry}
\usepackage{array}
\usepackage{longtable}
\usepackage{booktabs}
\usepackage{xcolor}
\usepackage{hyperref}

\geometry{margin=2.5cm}

\title{Relat\'orio do Dashboard de Monitoramento do TP02}
\author{Sistemas Distribu\'idos}
\date{\today}

\begin{document}

\maketitle

\section{Objetivo do Dashboard}

O dashboard foi criado para facilitar a visualiza\c{c}\~ao do funcionamento do protocolo
de exclus\~ao m\'utua distribu\'ida implementado no trabalho pr\'atico. A interface permite
acompanhar, de forma resumida, se os clientes est\~ao realizando requisi\c{c}\~oes, se os
peers do Cluster Sync est\~ao entrando corretamente na se\c{c}\~ao cr\'itica e se o recurso
compartilhado \texttt{R} est\'a sendo acessado sem sobreposi\c{c}\~oes.

As informa\c{c}\~oes exibidas s\~ao extra\'idas principalmente de dois arquivos:

\begin{itemize}
    \item \texttt{recurso\_R.txt}: arquivo que registra os eventos de entrada, sa\'ida e
    confirma\c{c}\~ao de acesso ao recurso compartilhado.
    \item \texttt{logs\_execucao.txt}: arquivo que registra as mensagens emitidas pelos
    clientes e pelos peers durante a execu\c{c}\~ao do sistema.
\end{itemize}

\section{Status Geral}

O campo de status geral indica uma conclus\~ao autom\'atica sobre a execu\c{c}\~ao observada.
Ele pode aparecer como funcionamento correto, aguardando eventos ou necessidade de
verifica\c{c}\~ao.

\begin{longtable}{>{\raggedright\arraybackslash}p{4cm} >{\raggedright\arraybackslash}p{10cm}}
\toprule
\textbf{Informa\c{c}\~ao} & \textbf{Significado} \\
\midrule
\texttt{Funcionamento correto} &
Indica que o dashboard encontrou eventos no recurso, que os contadores principais est\~ao
balanceados, que n\~ao houve sobreposi\c{c}\~ao de se\c{c}\~ao cr\'itica, que n\~ao faltam
timestamps de cliente e que n\~ao foram encontrados erros nos logs. \\
\texttt{Aguardando eventos} &
Indica que ainda n\~ao h\'a eventos suficientes em \texttt{recurso\_R.txt}. Isso pode
acontecer antes do in\'icio da execu\c{c}\~ao ou logo ap\'os limpar o arquivo. \\
\texttt{Verificar alertas} &
Indica que alguma valida\c{c}\~ao falhou, por exemplo: contadores diferentes, erro nos logs,
timestamp ausente, pedido incompleto ou poss\'ivel sobreposi\c{c}\~ao. \\
\bottomrule
\end{longtable}

\section{Resumo dos Eventos do Recurso}

Esta se\c{c}\~ao apresenta os principais contadores derivados do arquivo
\texttt{recurso\_R.txt}. Cada opera\c{c}\~ao correta de escrita no recurso deve gerar tr\^es
linhas: \texttt{ENTER\_CS}, \texttt{EXIT\_CS} e \texttt{COMMITTED}.

\begin{longtable}{>{\raggedright\arraybackslash}p{4cm} >{\raggedright\arraybackslash}p{10cm}}
\toprule
\textbf{Campo} & \textbf{Significado} \\
\midrule
\texttt{ENTER\_CS} &
Quantidade de vezes em que algum peer entrou na se\c{c}\~ao cr\'itica para acessar o recurso
compartilhado \texttt{R}. Cada entrada representa o in\'icio de uma opera\c{c}\~ao de escrita. \\
\texttt{EXIT\_CS} &
Quantidade de vezes em que algum peer saiu da se\c{c}\~ao cr\'itica. Para uma execu\c{c}\~ao
correta, esse valor deve ser igual ao total de \texttt{ENTER\_CS}. \\
\texttt{COMMITTED} &
Quantidade de opera\c{c}\~oes conclu\'idas e confirmadas. Para uma execu\c{c}\~ao correta,
esse valor tamb\'em deve ser igual ao total de \texttt{ENTER\_CS} e \texttt{EXIT\_CS}. \\
\texttt{Sobreposi\c{c}\~oes} &
Quantidade de casos em que um novo \texttt{ENTER\_CS} apareceu antes do \texttt{EXIT\_CS}
da opera\c{c}\~ao anterior. O valor esperado \'e zero, pois qualquer valor maior que zero
indica viola\c{c}\~ao da exclus\~ao m\'utua. \\
\texttt{Client timestamp} &
Quantidade de eventos do recurso que n\~ao possuem o timestamp enviado pelo cliente. O valor
esperado \'e zero, pois cada pedido do cliente deve carregar um timestamp para identificar
unicamente a solicita\c{c}\~ao. \\
\texttt{Erros nos logs} &
Quantidade de ocorr\^encias de termos associados a falhas em \texttt{logs\_execucao.txt},
como \texttt{Erro}, \texttt{ERROR}, \texttt{Traceback}, \texttt{timeout} ou
\texttt{Exception}. O valor esperado \'e zero. \\
\bottomrule
\end{longtable}

\section{Valida\c{c}\~oes Autom\'aticas}

O dashboard executa algumas verifica\c{c}\~oes autom\'aticas para avaliar se a execu\c{c}\~ao
observada respeita o protocolo.

\begin{longtable}{>{\raggedright\arraybackslash}p{5cm} >{\raggedright\arraybackslash}p{9cm}}
\toprule
\textbf{Valida\c{c}\~ao} & \textbf{Crit\'erio utilizado} \\
\midrule
Contadores balanceados &
Verifica se as quantidades de \texttt{ENTER\_CS}, \texttt{EXIT\_CS} e \texttt{COMMITTED}
s\~ao iguais. Se forem diferentes, alguma opera\c{c}\~ao ficou incompleta ou foi registrada
incorretamente. \\
Aus\^encia de sobreposi\c{c}\~ao &
Verifica se nunca houve dois peers simultaneamente dentro da se\c{c}\~ao cr\'itica. Essa
\'e a principal evid\^encia de que a exclus\~ao m\'utua foi respeitada. \\
Sequ\^encia correta de eventos &
Verifica se todo \texttt{EXIT\_CS} corresponde a um \texttt{ENTER\_CS} anterior e se n\~ao
existe se\c{c}\~ao cr\'itica aberta ao final do arquivo. \\
Presen\c{c}a de \texttt{client\_timestamp} &
Verifica se os eventos registrados em \texttt{recurso\_R.txt} possuem o timestamp original
enviado pelo cliente. \\
Pedidos completos &
Verifica se todo pedido que entrou na se\c{c}\~ao cr\'itica tamb\'em possui sa\'ida
\texttt{EXIT\_CS} e confirma\c{c}\~ao \texttt{COMMITTED}. \\
Aus\^encia de erros nos logs &
Verifica se o arquivo de logs n\~ao cont\'em mensagens indicativas de erro, exce\c{c}\~ao ou
timeout. \\
\bottomrule
\end{longtable}

\section{Pedidos por Cliente}

O gr\'afico de pedidos por cliente mostra quantas vezes cada cliente conseguiu acessar o
recurso compartilhado. A contagem \'e baseada nos eventos \texttt{ENTER\_CS}, pois cada
entrada na se\c{c}\~ao cr\'itica representa um pedido que come\c{c}ou a ser executado.

\begin{longtable}{>{\raggedright\arraybackslash}p{4cm} >{\raggedright\arraybackslash}p{10cm}}
\toprule
\textbf{Informa\c{c}\~ao} & \textbf{Significado} \\
\midrule
\texttt{client1}, \texttt{client2}, etc. &
Identifica o cliente que realizou pedidos de escrita ao recurso \texttt{R}. \\
Valor da barra &
Quantidade de acessos iniciados por aquele cliente. Segundo o enunciado, cada cliente deve
realizar entre 10 e 50 pedidos. \\
Utilidade &
Permite verificar se todos os clientes participaram da execu\c{c}\~ao e se a quantidade de
pedidos ficou dentro da faixa exigida. \\
\bottomrule
\end{longtable}

\section{Acessos por Peer}

O gr\'afico de acessos por peer mostra quantas vezes cada elemento do Cluster Sync entrou na
se\c{c}\~ao cr\'itica. No projeto, cada cliente conhece apenas um peer, ent\~ao a distribui\c{c}\~ao
por peer tende a acompanhar a distribui\c{c}\~ao por cliente.

\begin{longtable}{>{\raggedright\arraybackslash}p{4cm} >{\raggedright\arraybackslash}p{10cm}}
\toprule
\textbf{Informa\c{c}\~ao} & \textbf{Significado} \\
\midrule
\texttt{peer1}, \texttt{peer2}, etc. &
Identifica o processo do Cluster Sync que coordenou determinado acesso ao recurso
compartilhado. \\
Valor da barra &
Quantidade de vezes em que aquele peer entrou na se\c{c}\~ao cr\'itica. \\
Utilidade &
Permite verificar se os cinco processos do Cluster Sync participaram da execu\c{c}\~ao e se
o acesso ao recurso passou por diferentes peers. \\
\bottomrule
\end{longtable}

\section{Logs do Protocolo}

Esta se\c{c}\~ao resume informa\c{c}\~oes extra\'idas de \texttt{logs\_execucao.txt}. Ela ajuda
a conferir se o comportamento dos clientes e dos peers est\'a compat\'ivel com o enunciado.

\begin{longtable}{>{\raggedright\arraybackslash}p{5cm} >{\raggedright\arraybackslash}p{9cm}}
\toprule
\textbf{Campo} & \textbf{Significado} \\
\midrule
\texttt{COMMITTED clientes} &
Quantidade de mensagens \texttt{COMMITTED} recebidas pelos clientes. Essa mensagem indica
que o acesso ao recurso foi conclu\'ido e que o cliente pode iniciar seu per\'iodo de espera. \\
\texttt{Clientes finalizados} &
Quantidade de clientes que registraram a mensagem de finaliza\c{c}\~ao. Em uma execu\c{c}\~ao
completa do projeto, o valor esperado \'e 5. \\
\texttt{Sleeps} &
Quantidade de esperas realizadas pelos clientes ap\'os receberem \texttt{COMMITTED}. Cada
espera deve ocorrer antes de um novo pedido de escrita. \\
\texttt{Intervalo dos sleeps} &
Mostra o menor e o maior tempo de espera registrados. Pelo enunciado, os valores devem
estar entre 1 e 5 segundos. \\
\texttt{Sucessor imediato} &
Quantidade de vezes em que um peer liberou um \texttt{OK} para o pedido adiado com timestamp
imediatamente superior ao seu. Esse campo demonstra a regra de libera\c{c}\~ao para o sucessor
temporal ap\'os a sa\'ida da se\c{c}\~ao cr\'itica. \\
\texttt{Adiados seguintes} &
Quantidade de libera\c{c}\~oes de \texttt{OK} para outros pedidos adiados que vinham depois
do sucessor temporal imediato. Isso preserva o funcionamento do protocolo sem deixar peers
travados aguardando permiss\~oes antigas. \\
\texttt{Linhas de log} &
Quantidade total de linhas presentes em \texttt{logs\_execucao.txt}. Esse valor ajuda a
confirmar se o arquivo de logs foi preenchido durante a execu\c{c}\~ao. \\
\bottomrule
\end{longtable}

\section{\'Ultimos Eventos do Recurso R}

A tabela de \'ultimos eventos mostra as linhas mais recentes interpretadas a partir de
\texttt{recurso\_R.txt}. Ela permite acompanhar a execu\c{c}\~ao quase em tempo real.

\begin{longtable}{>{\raggedright\arraybackslash}p{4cm} >{\raggedright\arraybackslash}p{10cm}}
\toprule
\textbf{Coluna} & \textbf{Significado} \\
\midrule
\texttt{Evento} &
Tipo do evento registrado. Pode ser \texttt{ENTER\_CS}, \texttt{EXIT\_CS} ou
\texttt{COMMITTED}. \\
\texttt{Peer} &
Peer do Cluster Sync respons\'avel pelo acesso ao recurso naquele evento. \\
\texttt{Cliente} &
Cliente que originou o pedido de acesso ao recurso. \\
\texttt{Pedido} &
Identificador \'unico do pedido, por exemplo \texttt{client1-3}. \\
\texttt{Timestamp cliente} &
Timestamp gerado pelo cliente no momento em que ele enviou o pedido ao peer conhecido. Esse
campo evidencia que o pedido do cliente possui uma marca temporal pr\'opria. \\
\texttt{Timestamp evento} &
Timestamp registrado pelo peer no momento do evento no recurso. Ele permite observar a ordem
temporal dos acessos e medir o tempo dentro da se\c{c}\~ao cr\'itica. \\
\bottomrule
\end{longtable}

\section{Como Interpretar uma Execu\c{c}\~ao Correta}

Uma execu\c{c}\~ao pode ser considerada correta quando as seguintes condi\c{c}\~oes s\~ao
observadas no dashboard:

\begin{itemize}
    \item \texttt{ENTER\_CS}, \texttt{EXIT\_CS} e \texttt{COMMITTED} possuem o mesmo valor.
    \item O n\'umero de sobreposi\c{c}\~oes \'e zero.
    \item O n\'umero de timestamps de cliente ausentes \'e zero.
    \item O n\'umero de erros nos logs \'e zero.
    \item Os cinco clientes aparecem como finalizados ao final da execu\c{c}\~ao.
    \item Cada cliente possui entre 10 e 50 pedidos.
    \item Os sleeps dos clientes ficam entre 1 e 5 segundos.
    \item As dura\c{c}\~oes de se\c{c}\~ao cr\'itica ficam entre 0,2 e 1 segundo.
    \item Existem registros de libera\c{c}\~ao para o sucessor temporal imediato quando h\'a
    pedidos adiados.
\end{itemize}

\section{Conclus\~ao}

O dashboard funciona como uma ferramenta de apoio para demonstrar o comportamento do
protocolo de exclus\~ao m\'utua distribu\'ida. Ele n\~ao substitui a an\'alise do c\'odigo, mas
ajuda a apresentar evid\^encias pr\'aticas de que o sistema respeita a exclus\~ao m\'utua,
processa pedidos dos cinco clientes, registra timestamps e conclui as opera\c{c}\~oes com
\texttt{COMMITTED}.

\end{document}
